\documentclass[10pt]{article}
\usepackage[margin=0.55in]{geometry}
\usepackage{graphicx}
\usepackage{caption}
\usepackage{amsmath}
\usepackage{tikz}
\captionsetup{font=small,labelfont=bf}
\pagestyle{empty}
\begin{document}
\renewcommand{\thefigure}{1 amendment}
\begin{figure}[p]
\centering
\includegraphics[width=0.82\textwidth]{fig1_endpoint_normalization_amendment.pdf}
\caption{\textbf{Endpoint coordinates and canonical metrics.} The travel
optimum maps to $(0,1)$ and the cell-state optimum to $(1,0)$. Travel cost
$T$ and cell-state cost $S$ are normalized by the spans between these two
attained endpoint assignments. The natural lineage $L$ and the first-cousin
null mean $N$ use the same anchors. $P^*$ is the closest attained sampled-front
assignment to $L$ in these coordinates; $d_{LP}$ and $d_{NP}$ are measured
from $L$ and $N$ to this same assignment. Canonical position $u$ is the
fraction of front arc length from the travel optimum to $P^*$. Each subtree
uses its own endpoints for shape summaries, while assignment comparisons
within a cost space share one reference. This display transformation leaves
the optimization and selected assignments unchanged. This independent
component awaits integration into the co-author's Figure~1.}
\label{fig:ce_endpoint_normalization_amendment}
\end{figure}
\end{document}
